\section{Introduction}
\label{sec:intro}

Vision-Language Models (\vlm{}s) have become standard input components in everything from hosted assistants to tool-using agents to MCP-style retrieval pipelines. Recent universal-attack work has shown that an adversarially-perturbed image, optimised once and reused across queries, can override the safety alignment of small open \vlm{}s with reported attack success rates in the $60$--$80\%$ range \citep{rahmatullaev2025universal,qi2024visual,schlarmann2023robustness,bagdasaryan2023imagessounds}. Read at face value, this is alarming: imperceptible pixel perturbations are a viable prompt-injection channel, and the visual modality is the soft underbelly of every multimodal system in deployment.

We argue that the picture is less dramatic, and more interesting, than the headline number suggests. Reported ASR conflates two distinct events:

\begin{itemize}
  \item \textbf{Output Affected (drift).} The \vlm{}'s response was perturbed in some way --- it became different from what it would have said on the clean image.
  \item \textbf{Target Injected (payload delivery).} The attacker's chosen target string actually appeared in the response.
\end{itemize}

These are different things. A response that flips from ``a dog in a field'' to ``a collage of incoherent text fragments'' is affected but contains no payload. A response that emits the target URL inside an otherwise-on-topic answer is a different kind of failure. Universal-attack benchmarks that score success via an LLM-judge prompt of the form ``does the answer differ from the clean baseline'' systematically conflate the two and overstate the attack's value as a delivery vector.

\paragraph{This paper.}
We compose two existing universal-attack methods --- the Universal Adversarial Attack of \citet{rahmatullaev2025universal} (Stage 1) and the \textsc{AnyAttack} self-supervised encoder--decoder of \citet{zhang2025anyattack} (Stage 2) --- under an \Linf{} budget of $\eps = 16/255$ and PSNR $\approx 25.2$~dB. We add a dual-axis evaluation that scores Influence (Ratcliff-Obershelp drift, deterministic baseline; LLM-judged 4-tier ordinal level) and Precise Injection (LLM-judged 4-tier categorical: \texttt{none / weak / partial / confirmed}) independently, using \textsc{DeepSeek-V4-Pro} \citep{deepseekv3} as the LLM judge in thinking mode (calibrated against Claude Opus 4.7 \citep{anthropic2025opus47} with Cohen's $\kappa = 0.77$ on the injection axis, substantial agreement \citep{landiskoch1977}). We ship the entire LLM-call cache so reviewers can re-derive paper numbers bit-exact without an API key. We run the resulting pipeline on a $7 \times 7 \times 3$ matrix --- seven attack prompts, seven clean test images, three white-box surrogate ensembles --- against four open \vlm{}s, producing $6{,}615$ (clean, adversarial) response pairs.

The headline finding is the gap. Across the same $6{,}615$ pairs, programmatic disruption sits at $66.4\%$ and LLM-judged disruption (substantial+complete) at $46.6\%$, but Precise Injection sits at $0.756\%$ broad / $0.030\%$ verbatim --- a roughly $90\times$ divergence even at the broadest definition of injection. Disruption is broad and architecture-dependent (every transformer-style \vlm{} we test is disturbed on $\geq 99\%$ of pairs by the deterministic measure); payload delivery is rare and image-dependent. The few literal injections that do land cluster on screenshot-like images whose semantics already contain URL- or account-shaped tokens. BLIP-2 shows \emph{zero detectable drift} at $\Linf = 16/255$ across all $2{,}205$ pairs even when used as a Stage-1 surrogate.

\paragraph{Contributions.}
This paper makes four contributions:
\begin{itemize}
  \item \textbf{C1.} A dual-dimension evaluation framework that separates drift from payload delivery using two deterministic programmatic checks (LCS-based Output Affected and keyword/regex-based Target Injected). The evaluation runs end-to-end in $\sim 5$ minutes on a CPU and removes the LLM-judge from the scoring loop. (\S\ref{sec:method-stage3})
  \item \textbf{C2.} A $6{,}615$-pair systematic sweep over four open \vlm{}s, seven prompts, and seven test images, showing the dual-axis split. The two axes diverge by roughly $90\times$ on the same data ($66.4\%$ programmatic disruption vs $0.756\%$ broad injection by the v3 LLM judge), demonstrating that single-number ``ASR'' metrics conflate two qualitatively different failure modes. (\S\ref{sec:results})
  \item \textbf{C3.} An architectural observation: BLIP-2 shows \emph{zero detectable drift} at $\Linf = 16/255$ ($0/2{,}205$ pairs registered any Output-Affected score $> 0$) even when included as a Stage-1 surrogate. We enumerate three candidate causes (Stage-2 decoder fusion, the $448\!\to\!224$ resolution + Q-Former double bottleneck, and gradient dilution at training time) and propose a one-shot ablation that would distinguish them. We do not claim a definitive mechanism. (\S\ref{sec:discussion})
  \item \textbf{C4.} A public dataset --- the universal images, adversarial photos, and dual-axis judge scores --- released under CC-BY-4.0 at \texttt{huggingface.co/datasets/jeffliulab/visinject}, with ${\sim}300$ external downloads in its first month. (\S\ref{sec:conclusion})
\end{itemize}

\paragraph{Roadmap.}
\S\ref{sec:related} positions this work against the universal-attack, prompt-injection, and VLM-jailbreak literatures. \S\ref{sec:threat} states the threat model. \S\ref{sec:method} describes the three-stage pipeline. \S\ref{sec:experiments} lays out the experimental setup. \S\ref{sec:results} reports aggregate results, \S\ref{sec:cases} walks through three case studies, and \S\ref{sec:discussion} discusses why drift outpaces injection and what BLIP-2's immunity means for VLM defense design. \S\ref{sec:conclusion} closes with the released artifacts.
